\documentclass[11pt,a4paper]{article}

\usepackage[utf8]{inputenc}
\usepackage[T1]{fontenc}
\usepackage[polish]{babel}
\usepackage{lmodern}
\usepackage[margin=2.5cm]{geometry}
\usepackage{amsmath,amssymb}
\usepackage{graphicx}
\usepackage{booktabs}
\usepackage{enumitem}
\usepackage{hyperref}
\usepackage{url}
\usepackage{listings}
\usepackage{xcolor}

\lstset{
    basicstyle=\ttfamily\small,
    breaklines=true,
    columns=fullflexible,
    backgroundcolor=\color{gray!10},
    frame=single,
}

\title{Raport końcowy \\
\large Zastosowanie Monte Carlo Tree Search (MCTS/UCT) do gry w warcaby 10$\times$10 bez obowiązkowego bicia z damkami bijącymi po całej diagonali}
\author{Jakub Kindracki \and Wiktor Kobielski}
\date{\today}

\begin{document}
\maketitle

\begin{abstract}
Niniejszy dokument jest szkieletem (template) raportu końcowego. W obecnej
wersji opisuje wyłącznie \emph{plan i metodykę testów} dla hipotez H1--H4 --
konfiguracje eksperymentów, sposób ich uruchamiania (pipeline
runy~$\to$~statystyki~$\to$~testy istotności~$\to$~wykresy), kryteria
potwierdzenia hipotez oraz zapewnienie powtarzalności. Sekcje z wynikami
(tabele, wykresy, dyskusja) zostaną uzupełnione po wykonaniu eksperymentów.
\end{abstract}

\tableofcontents

\section{Wprowadzenie}

\subsection{Zakres tego dokumentu}

Ten plik (\texttt{documentation/final/raport.tex}) jest \emph{szablonem}
raportu końcowego. Sekcja~\ref{sec:metodologia} opisuje wspólną
infrastrukturę eksperymentalną, a sekcje~\ref{sec:h1}--\ref{sec:h4}
opisują -- dla każdej z czterech hipotez badawczych -- dokładny plan testów:
jakie konfiguracje graczy są porównywane, ile partii zostanie rozegranych,
jakie kryterium decyduje o potwierdzeniu hipotezy oraz jakie artefakty
(pliki JSON/CSV/PNG) powstaną. Wyniki (sekcje ,,Wyniki'' w każdej
podsekcji) są obecnie placeholderami -- zostaną wypełnione po uruchomieniu
skryptów \texttt{run\_h1.py}--\texttt{run\_h4.py}.

\subsection{Przypomnienie hipotez}

\begin{description}[leftmargin=2.2cm,style=nextline]
    \item[H1 (UCT vs heurystyka).] Czysty UCT z wystarczającym limitem
        iteracji ($\geq 10\,000$ na ruch) gra istotnie skuteczniej od
        gracza czysto heurystycznego, osiągając $> 60\%$ wygranych w
        meczu z naprzemienną zmianą kolorów.
    \item[H2 (RAVE poprawia UCT).] UCT z rozszerzeniem RAVE wygrywa
        $> 60\%$ partii z czystym UCT przy tym samym limicie iteracji.
    \item[H3 (Progressive bias poprawia UCT).] UCT z heurystycznym
        progressive bias wygrywa $> 60\%$ partii z czystym UCT, a jego
        przewaga jest największa przy małym limicie iteracji
        ($\leq 2\,000$/ruch); wraz ze wzrostem liczby iteracji oczekujemy
        zaniku przewagi.
    \item[H4 (krawędzie vs środek planszy).] Heurystyki premiujące figury
        przy krawędzi planszy są skuteczniejsze niż heurystyki premiujące
        figury w centrum -- weryfikacja w dwóch konfiguracjach
        (gracz heurystyczny oraz UCT + progressive bias), oba z wariantem
        ,,krawędzie'' vs ,,środek''. Hipoteza potwierdzona, jeśli wariant
        ,,krawędzie'' wygrywa $> 50\%$ partii w \emph{obu} konfiguracjach.
\end{description}

\section{Metodologia i infrastruktura wspólna}
\label{sec:metodologia}

\subsection{Pipeline eksperymentalny}

Każda hipoteza ($H_i$, $i=1,\dots,4$) jest weryfikowana przez jeden
samodzielny skrypt \texttt{run\_h\{i\}.py}, który realizuje cały
pipeline w jednym przebiegu:

\begin{enumerate}
    \item \textbf{Uruchomienia} (\emph{runs}) -- rozegranie partii
        AI vs AI dla wszystkich konfiguracji danej hipotezy, zapis
        surowych wyników do \texttt{output/h\{i\}/raw\_results.json}
        oraz replayable logów partii do \texttt{output/h\{i\}/games/}
        (\texttt{hypothesis\_lib.run\_stage1}).
    \item \textbf{Statystyki} -- dla każdej konfiguracji: łączny
        odsetek wygranych $p_1$ (wraz z 95\% przedziałem ufności
        Wilsona), liczba partii, remisów itd.
        (\texttt{analyze\_results.compute\_stats}).
    \item \textbf{Testy hipotez i istotności} -- dwustronny dokładny test
        dwumianowy przeciwko $p_0=0{,}5$ (na partiach zdecydowanych,
        bez remisów) oraz test spójności między ziarnami (czy $p_1$
        wygrywa w większości ziaren bazowych). Hipoteza jest uznawana za
        potwierdzoną, gdy: $p_1\text{-win-rate} > \text{próg}$ \textbf{i}
        $p\text{-value} < 0{,}05$ \textbf{i} wynik jest spójny między
        ziarnami (zob. \texttt{analyze\_results.analyze\_config}). Próg
        wynosi $60\%$ dla H1--H3 oraz $50\%$ dla H4 (zgodnie z
        formułowaniem hipotez w konspekcie).
    \item \textbf{Wykresy} -- wykres słupkowy odsetka wygranych z
        przedziałami ufności (oraz liniami referencyjnymi $0{,}5$ i
        progu) dla każdej konfiguracji, a dla hipotez ze sweepem
        budżetu iteracji ($H_1$--$H_3$) dodatkowo wykres
        odsetek-wygranych-vs-liczba-iteracji z pasmem CI
        (\texttt{plot\_results.generate\_plots}). Pliki PNG trafiają do
        \texttt{output/h\{i\}/plots/}.
\end{enumerate}

Wszystkie cztery etapy wykonują się automatycznie w ramach jednego
wywołania, tj. \texttt{python run\_h\{i\}.py} bez żadnych dodatkowych
flag produkuje kompletny zestaw artefaktów: \texttt{raw\_results.json},
\texttt{stats.json}, \texttt{stats.csv} oraz pliki PNG w
\texttt{plots/}.

\subsection{Powtarzalność (reprodukowalność)}
\label{sec:reprodukowalnosc}

Zgodnie z konspektem każda konfiguracja eksperymentu jest powtarzana dla
$\geq 5$ różnych ziaren bazowych
(\texttt{hypothesis\_lib.derive\_seeds}, domyślnie ziarna bazowe
$0, 10007, 20014, 30021, 40028$ -- pochodne ziarna pochodne ziarna
bazowego $0$). Każdy gracz przyjmuje parametr \texttt{seed}; ziarna
poszczególnych partii i ruchów są wyprowadzane deterministycznie z ziarna
bazowego (\texttt{experiments.run\_match}).

Root-parallel MCTS dzieli budżet iteracji między procesy robocze i
przypisuje każdemu procesowi ziarno zależne od jego indeksu, więc liczba
workerów jest częścią tego, co determinuje wynik wyszukiwania. Aby
domyślne (bezargumentowe) wywołanie \texttt{python run\_h\{i\}.py} było w
pełni powtarzalne niezależnie od liczby rdzeni dostępnej maszyny, domyślna
liczba workerów jest ustalonym stałą (\texttt{hypothesis\_lib.\allowbreak
DEFAULT\_WORKERS = 4}), a nie \texttt{os.cpu\_count()}. Flaga
\texttt{--workers N} pozwala przyspieszyć (i ,,zbogacić'') eksperyment na
maszynach wieloprocesorowych, kosztem utraty ścisłej powtarzalności
między maszynami o różnej liczbie rdzeni.

Każda rozegrana partia jest zapisywana w formacie tekstowym (sekwencja
ruchów + metadane: gracze, heurystyki, ziarna, liczba odwiedzonych węzłów
drzewa) i może zostać odtworzona za pomocą
\texttt{experiments.replay\_game\_log}.

\subsection{Budżet czasowy}

Każdy skrypt \texttt{run\_h\{i\}.py} ma domyślny miękki budżet czasowy
\texttt{--time-budget-min~110} (czyli ok. 1h50min), tak aby cały
pipeline (uruchomienia + statystyki + wykresy) zmieścił się w limicie
2 godzin na hipotezę. Liczba partii na ziarno jest dynamicznie
dopasowywana do tego budżetu na podstawie krótkiej kalibracji prędkości
MCTS na danej maszynie (\texttt{hypothesis\_lib.\allowbreak
calibrate\_seconds\_per\_iteration},
\texttt{allocate\_games\_per\_seed}) -- każda konfiguracja otrzymuje co
najmniej $5 \times 1 = 5$ partii (1 partia/ziarno $\times$ 5 ziaren), a
pozostały budżet jest zachłannie rozdzielany między konfiguracje
(do max. 10 partii/ziarno = 50 partii/konfigurację).

\subsection{Mierzone wielkości}

Dla każdej partii rejestrowane są: zwycięzca (lub remis), liczba ruchów,
czas trwania oraz średnia liczba odwiedzonych węzłów drzewa MCTS na ruch
(\texttt{last\_node\_count} / \texttt{mcts.\_count\_nodes}) -- miara
,,wysiłku obliczeniowego'' wyszukiwania, zwracana w
\texttt{stats.json}/\texttt{stats.csv} jako \texttt{p1\_avg\_nodes} /
\texttt{p2\_avg\_nodes} (w surowych wynikach) i wykorzystywana do
interpretacji wyników (np. czy przewaga UCT+RAVE/progressive bias wynika
z efektywniejszego przeszukiwania, a nie tylko z większego budżetu).


\section{H1: UCT vs heurystyka}
\label{sec:h1}

\subsection{Plan testów}

\textbf{Skrypt:} \texttt{run\_h1.py} (domyślny katalog wyjściowy:
\texttt{output/h1/}).

\textbf{Porównywani gracze:} \texttt{MCTSPlayer(variant="uct")} ($p_1$)
vs \texttt{HeuristicPlayer} ($p_2$, heurystyka \texttt{"base"}).

\textbf{Sweep budżetu iteracji ($p_1$):}
$\{500,\ 1\,000,\ 2\,000,\ 5\,000,\ 10\,000\}$ -- pięć konfiguracji
\texttt{uct500}, \dots, \texttt{uct10000}. Tylko $p_1$ wykonuje
przeszukiwanie MCTS (\texttt{mcts\_search\_units = iterations}), więc
koszt konfiguracji rośnie liniowo z liczbą iteracji. Wartość
$10\,000$ odpowiada literalnemu warunkowi H1 ($\geq 10\,000$/ruch).

\textbf{Liczba partii:} każda konfiguracja $\geq 5$ ziaren bazowych
$\times \geq 1$ partia/ziarno (kolory zamieniane między parami partii w
ramach \texttt{run\_match}); tańsze konfiguracje (np. \texttt{uct500})
otrzymają do 50 partii w ramach dostępnego budżetu czasowego, droższe
(\texttt{uct10000}) co najmniej 5 partii.

\textbf{Kryterium potwierdzenia:} dla danej konfiguracji H1 jest uznana za
potwierdzoną, gdy odsetek wygranych UCT $> 60\%$, $p$-value (dwustronny
test dwumianowy vs. $0{,}5$) $< 0{,}05$ i UCT wygrywa w większości z $\geq
5$ ziaren. Oczekiwany wynik: hipoteza potwierdzona przy $10\,000$
iteracjach (i prawdopodobnie przy mniejszych budżetach).

\textbf{Artefakty:}
\begin{itemize}
    \item \texttt{output/h1/raw\_results.json} -- surowe wyniki per
        konfiguracja $\times$ ziarno.
    \item \texttt{output/h1/games/*.txt} -- replayable logi partii.
    \item \texttt{output/h1/stats.json}, \texttt{stats.csv} -- statystyki
        i werdykty per konfiguracja.
    \item \texttt{output/h1/plots/h1\_winrate\_bar.png} -- odsetek
        wygranych UCT z 95\% CI per konfiguracja, z liniami
        referencyjnymi $0{,}5$ i $0{,}6$.
    \item \texttt{output/h1/plots/h1\_winrate\_vs\_iterations.png} --
        odsetek wygranych UCT vs. budżet iteracji (skala logarytmiczna),
        z pasmem 95\% CI.
\end{itemize}

\subsection{Wyniki}

\textit{[TODO: tabela \texttt{stats.csv} oraz wykresy z
\texttt{output/h1/plots/} po wykonaniu \texttt{python run\_h1.py}.]}


\section{H2: UCT + RAVE vs UCT}
\label{sec:h2}

\subsection{Plan testów}

\textbf{Skrypt:} \texttt{run\_h2.py} (domyślny katalog wyjściowy:
\texttt{output/h2/}).

\textbf{Porównywani gracze:} \texttt{MCTSPlayer(variant="rave")} ($p_1$)
vs \texttt{MCTSPlayer(variant="uct")} ($p_2$), \textbf{przy tym samym
budżecie iteracji} dla obu stron.

\textbf{Sweep budżetu iteracji (wspólny dla obu stron):}
$\{500,\ 1\,000,\ 2\,000,\ 5\,000\}$ -- cztery konfiguracje
\texttt{rave500}, \dots, \texttt{rave5000}. Obie strony wykonują
przeszukiwanie MCTS (\texttt{mcts\_search\_units = 2 * iterations}), więc
koszt konfiguracji jest ok. dwukrotnie wyższy niż analogicznej
konfiguracji H1 -- stąd mniejszy (4-elementowy) sweep, aby zmieścić się w
budżecie 2h.

\textbf{Liczba partii:} analogicznie do H1 -- $\geq 5$ ziaren $\times \geq
1$ partia/ziarno, do 50 partii na konfigurację w ramach budżetu czasowego.

\textbf{Kryterium potwierdzenia:} RAVE wygrywa $> 60\%$ partii z UCT przy
tym samym budżecie, $p < 0{,}05$, spójnie między ziarnami.

\textbf{Artefakty:} analogicznie do H1 --
\texttt{output/h2/\{raw\_results.json, stats.json, stats.csv\}},
\texttt{output/h2/games/*.txt},
\texttt{output/h2/plots/\{h2\_winrate\_bar.png,
h2\_winrate\_vs\_iterations.png\}}.

\subsection{Wyniki}

\textit{[TODO: tabela \texttt{stats.csv} oraz wykresy z
\texttt{output/h2/plots/} po wykonaniu \texttt{python run\_h2.py}.]}


\section{H3: UCT + progressive bias vs UCT}
\label{sec:h3}

\subsection{Plan testów}

\textbf{Skrypt:} \texttt{run\_h3.py} (domyślny katalog wyjściowy:
\texttt{output/h3/}).

\textbf{Porównywani gracze:} \texttt{MCTSPlayer(variant="progressive")}
($p_1$, heurystyka \texttt{"base"} jako składnik biasu $h(\text{stan})$)
vs \texttt{MCTSPlayer(variant="uct")} ($p_2$), przy tym samym budżecie
iteracji.

\textbf{Sweep budżetu iteracji (wspólny dla obu stron):}
$\{500,\ 1\,000,\ 2\,000,\ 5\,000,\ 10\,000\}$ -- pięć konfiguracji
\texttt{progressive500}, \dots, \texttt{progressive10000}
(\texttt{mcts\_search\_units = 2 * iterations}, jak w H2). Sweep
obejmuje wartości $\leq 2\,000$ \emph{i} $> 2\,000$, aby na wykresie
odsetek-wygranych-vs-iteracje pokazać oczekiwany \emph{zanik przewagi}
progressive bias przy większych budżetach.

\textbf{Liczba partii:} analogicznie do H1/H2 -- $\geq 5$ ziaren $\times
\geq 1$ partia/ziarno, do 50 partii na konfigurację w ramach budżetu
czasowego.

\textbf{Kryterium potwierdzenia:} dla każdej konfiguracji -- progressive
bias wygrywa $> 60\%$ partii z UCT, $p < 0{,}05$, spójnie między ziarnami.
Dodatkowo (analiza opisowa, nie pojedynczy test statystyczny): czy odsetek
wygranych progressive bias maleje monotonicznie (lub przynajmniej
generalnie spada) ze wzrostem budżetu iteracji, z największą przewagą przy
$\leq 2\,000$ iteracjach -- widoczne na wykresie
\texttt{h3\_winrate\_vs\_iterations.png}.

\textbf{Artefakty:} analogicznie do H1 --
\texttt{output/h3/\{raw\_results.json, stats.json, stats.csv\}},
\texttt{output/h3/games/*.txt},
\texttt{output/h3/plots/\{h3\_winrate\_bar.png,
h3\_winrate\_vs\_iterations.png\}}.

\subsection{Wyniki}

\textit{[TODO: tabela \texttt{stats.csv} oraz wykresy z
\texttt{output/h3/plots/} po wykonaniu \texttt{python run\_h3.py}. W
szczególności: czy krzywa odsetka wygranych vs. iteracje jest malejąca, i
czy przewaga jest największa przy $\leq 2\,000$ iteracjach.]}


\section{H4: krawędzie vs środek planszy}
\label{sec:h4}

\subsection{Plan testów}

\textbf{Skrypt:} \texttt{run\_h4.py} (domyślny katalog wyjściowy:
\texttt{output/h4/}).

Hipoteza jest weryfikowana w \textbf{dwóch konfiguracjach}, w obu wariant
,,krawędzie'' (\texttt{heuristic="edge"}) jest $p_1$, a wariant ,,środek''
(\texttt{heuristic="center"}) jest $p_2$:

\begin{enumerate}
    \item \textbf{\texttt{heuristic\_edge\_vs\_center}} --
        \texttt{HeuristicPlayer(heuristic="edge")} vs
        \texttt{HeuristicPlayer(heuristic="center")}. Żadna ze stron nie
        wykonuje MCTS (\texttt{mcts\_search\_units = 0}), więc ta
        konfiguracja jest praktycznie darmowa -- otrzymuje maksymalną
        liczbę partii dopuszczalną przez \texttt{--max-games-per-seed}
        (domyślnie do 50).
    \item \textbf{\texttt{progressive\_edge\_vs\_center\_2000}} --
        \texttt{MCTSPlayer(variant="progressive", heuristic="edge")} vs
        \texttt{MCTSPlayer(variant="progressive", heuristic="center")},
        przy budżecie $2\,000$ iteracji (\texttt{mcts\_search\_units =
        2 * iterations = 4000}) -- wartość $2\,000$ leży w zakresie
        ,,małego budżetu'' z H3, gdzie wpływ biasu $h$ jest największy.
\end{enumerate}

\textbf{Liczba partii:} obie konfiguracje $\geq 5$ ziaren $\times \geq 1$
partia/ziarno; konfiguracja heurystyczna (1) skaluje się do limitu
\texttt{--max-games-per-seed} (domyślnie 50 partii), konfiguracja
progressive (2) -- analogicznie do H1--H3, w ramach pozostałego budżetu
czasowego.

\textbf{Kryterium potwierdzenia:} H4 jest uznana za potwierdzoną, gdy w
\emph{obu} konfiguracjach wariant ,,krawędzie'' wygrywa $> 50\%$ partii
(\texttt{stats.json: both\_configs\_hold}); per-konfiguracyjny
\texttt{hypothesis\_holds} dodatkowo wymaga $p < 0{,}05$ i spójności
między ziarnami (próg $50\%$, nie $60\%$, w odróżnieniu od H1--H3).

\textbf{Artefakty:}
\texttt{output/h4/\{raw\_results.json, stats.json, stats.csv\}},
\texttt{output/h4/games/*.txt},
\texttt{output/h4/plots/h4\_winrate\_bar.png} (wykres
odsetek-wygranych-vs-iteracje nie jest generowany dla H4, ponieważ
konfiguracje mieszają różne typy graczy i budżety iteracji).

\subsection{Wyniki}

\textit{[TODO: tabela \texttt{stats.csv} oraz wykres z
\texttt{output/h4/plots/} po wykonaniu \texttt{python run\_h4.py}. W
szczególności wartość \texttt{both\_configs\_hold} ze \texttt{stats.json}.]}


\section{Testy z udziałem ludzi}

\textit{[TODO: opis przebiegu testów z $\geq 5$ osobami, ankieta,
podsumowanie wyników -- po przeprowadzeniu testów z wykorzystaniem
interfejsu Flask (\texttt{app.py}).]}


\section{Podsumowanie}

\textit{[TODO: zbiorcze podsumowanie wyników H1--H4 oraz testów z
udziałem ludzi, wnioski, ograniczenia i propozycje dalszych badań.]}

\end{document}
